% !TeX root = ../ba-lentz.tex

\documentclass[ba-lentz.tex]{subfile}

\begin{document}

{\let\clearpage\relax \chapter{Geplante Vorgehensweise und Vorläufiger Zeitplan}}\label{summary_conclusion}

\section{Summary}\label{summary}

\section{Conclusion}\label{conclusion}

Diese Arbeit besteht aus 2 möglichen Schritten.
Zunächst sollen die Remote Sensing Daten, sowie die zu verbessernden dynamischen Vorhersagen für den Postprocessing Schritt vorbereitet werden.
Dazu sollen die Niederschlags und Temperaturdaten auf eine gemeinsame Auflösung gebracht und aus ihnen die $PET$ und schließlich die $KWB$ errechnet werden.
Anschließend sollen Masken erstellt werden, welche Regionen mit größerem und niedrigerem Vorhersageskill darstellen.
Mithilfe der Masken und den vorbereiteten Inputdaten soll dann das PCNN trainiert und mit dem trainierten PCNN die saisonalen $KWB$ Vorhersagen des DWD nachbearbeitet werden.
Zum Abschluss diesen ersten Teiles soll eine Evaluierung des Effekts, welchen die Nachbearbeitung auf die Vorhersage von Dürreereignissen hat, durchgeführt werden. \\
Falls diese Analyse im geplanten Zeitrahmen abgeschlossen werden kann, soll anschließend ein Downscaling Mechanismus in das PCNN eingebaut werden.
Anschließend soll die Analyse des ersten Teils wiederholt werden und mit einem anschließenden Downscaling Schritt, welcher die Vorhersagedaten auf die ursprüngliche Auflösung der Beobachtungsdaten hebt, supplementiert werden.
Schließlich sollen die Ergebnisse noch einmal ausführlich evaluiert und in den Kontext anderer Biaskorrekturen oder hybrider Vorhersagemethoden gesetzt werden. \\
Zeitlich soll die Vorbereitung der Daten in den ersten zwei Wochen nach Bearbeitungsstart abgeschlossen sein.
Das Anschließende Training und Evaluieren des PCNNs sollte weitere 3-4 Wochen in Anspruch nehmen.
Falls dies gelingt soll der zweite Teil der Analyse in weiteren 2-4 Wochen durchgeführt werden.
Für das Zusammentragen und Niederschreiben der Ergebnisse werden zum Schluss weitere 2-4 Wochen eingeplant.


\end{document}